\ssr{ВВЕДЕНИЕ}

Большинство объектов трёхмерной сцены описываются поверхностями. Для их изображения требуется определить, какие поверхности видны наблюдателю, как они освещены и какие тени отбрасывают, и для каждой из этих задач в компьютерной графике существуют классические алгоритмы. Северное сияние поверхности не имеет: это полупрозрачная светящаяся среда, яркость которой складывается из излучения её точек, лежащих на луче зрения, с учётом поглощения. Кроме того, форма и яркость занавесов сияния меняются во времени, а их свечение само освещает окружающую сцену. Поэтому для построения изображения северного сияния помимо классических алгоритмов требуются методы синтеза и визуализации светящейся среды.

Цель работы "--- разработка программного обеспечения для визуализации трёхмерного изображения северного сияния.

Для достижения поставленной цели требуется решить следующие задачи:
\begin{enumerate}
	\item формализовать объекты синтезируемой сцены и описать ограничения на входные данные;
	\item проанализировать методы и алгоритмы решения следующих задач и выбрать алгоритмы для реализации:
	\begin{itemize}
		\item расчёт освещённости поверхностей;
		\item построение теней;
		\item синтез формы занавеса сияния;
		\item визуализация светящейся среды;
		\item освещение сцены свечением сияния;
	\end{itemize}
	\item спроектировать программное обеспечение;
	\item выбрать средства реализации и реализовать спроектированное программное обеспечение;
	\item исследовать зависимость времени синтеза кадра от шага интегрирования и от размера расчётной сетки.
\end{enumerate}
